\subsection{Rootable actions and the same original Rees components}

\begin{lemma}[A rootable group fixes the original minimal primes]
\label{algebra:rootable-action-fixes-minimal-primes}
\lean{EnvelopeProperty.minimalPrime_comap_eq_self_of_rootable_action}
\leanok
Let $E$ be ANY Noetherian commutative ring and $G$ a group for which
every positive integer power map is surjective. Let
$\rho:G\to\operatorname{Aut}_{\mathrm{ring}}(E)$ be a group
homomorphism, and let $I\subseteq E$ be an original ideal satisfying
$\rho(s)^{-1}(I)=I$ for every $s\in G$.
Then EVERY original minimal prime $P$ over $I$ satisfies
\[
 \rho(s)^{-1}(P)=P\qquad(s\in G).
\]
There is no domain, reducedness, field, commutativity of $G$, finite
normalization, or fixed-prime hypothesis.
\end{lemma}
\begin{proof}
Let $M$ be the set of original minimal primes over $I$. It is finite
by Noetherianity. Surjectivity of each original ring automorphism and
the supplied equality for $I$ show that inverse image maps $M$ into
itself. Define
\[
 \pi(s)(P)=\rho(s^{-1})^{-1}(P).
\]
The inverse permutation is $\pi(s^{-1})$. The inverse in this formula
is essential: inverse image reverses composition, so the group and
automorphism laws give $\pi(st)=\pi(s)\pi(t)$ with the ordinary
composition law on permutations. Thus $\pi$ is a genuine group
homomorphism to the finite group $\operatorname{Perm}(M)$.

Set $N=|\operatorname{Perm}(M)|$. This is a positive integer even
when $M$ is empty. Given $s\in G$, choose an actual $z\in G$ with
$z^N=s$. The existing finite-group power theorem gives
$\pi(s)=\pi(z)^N=1$. Therefore $\pi$ is trivial. Apply this to
$s^{-1}$ and the original $P$ to obtain exactly the asserted inverse
image equality. No representative of a minimal prime or quotient of
the original ring is substituted.
\end{proof}

\begin{theorem}[Every original multigraded Rees component has no smaller order]
\label{algebra:actual-multigraded-rees-component-divisibility}
\leanok
\uses{def:gl-model,def:bj-rees,algebra:rootable-action-fixes-minimal-primes,lem:actual-multigraded-torus-component-extraction,lem:actual-original-rees-component-inputs}
Let $k$ be algebraically closed and let $X$ be an ORIGINAL integral
multiprojective model over $k$, with original coefficient field $K$
and original finitely generated multigraded subalgebra
$S\subseteq K[Z_1,\ldots,Z_r]$.
Write $\operatorname{homComp}_a(p)$ for its existing original component
of multidegree $a\in\mathbb N^r$, as an element of the SAME $S$.
Retain $B=S[t]$ and its SAME original fraction field $L$.

Let $\iota$ be any finite index type. For each $i$, take actual
$a_i\in\mathbb N^r$, $n_i\in\mathbb N$, and $f_i\in K$ with
$f_i\in X.\operatorname{piece}(a_i)$, and define the actual original family
\[
 g_i=t^{n_i}\,X.\operatorname{ofPiece}(f_i)\in B.
\]
For EVERY actual original Rees index $x$ of this family, put
$V_x=\operatorname{reesValuationRing}(L,g,x)$ and retain the original
inclusion $\eta_x:B\to V_x$.
Then EVERY $F\in B$, $n\in\mathbb N$ and $a\in\mathbb N^r$ satisfy
\[
 \eta_x(F)\mid
 \eta_x\bigl(t^n\operatorname{homComp}_a(F_n)\bigr),
 \qquad F_n=\operatorname{coeff}_n(F).
 \tag{1}
\]
Equivalently the original normalized DVR order of every such joint
component is at least that of $F$. Neither scale stability of $V_x$,
component divisibility, a competing-point family, nor a separator is
an input. No normality of $S$, projective section-surjectivity,
positivity of the order of $t$, or generation of a unit ideal is assumed.
\end{theorem}
\begin{proof}
The original model makes $S$ a domain and a finite-type $k$-algebra.
The SAME is true of $B=S[t]$. Algebraic closedness supplies perfectness,
so the existing original finite-normalization theorem makes every
normalized Rees chart Noetherian. Its original minimal-prime
localizations are therefore the existing genuine DVRs. All original
coefficient actions and fraction-field inclusions are retained.

Fix one coordinate block $c$ and $u\in k^\times$. Use the EXISTING
ambient endomorphism $Z_c\mapsto uZ_c$, fixing all other variables
and every coefficient in $K$. It preserves the original $S$:
this is the existing scale-stability of the original image subalgebra,
with coefficient extension $K\to K$ and its existing equality with
$S$. Scaling by $u^{-1}$ is its inverse, as is checked on each original
coefficient. The two restrictions are genuine inverse automorphisms of
the SAME $S$, preserve its original $k$ action, and satisfy the group
law on every original coefficient. Apply the existing polynomial
map-equivalence to the coefficients of $B$, fixing $t$. Denote the
resulting automorphism by $\sigma_c(u)$. The existing monomial scaling
formula gives the WHOLE original identity
\[
 \sigma_c(u)(g_i)=u^{a_i(c)}g_i.
 \tag{2}
\]
These are units coming from the original $k$ action.

Extend $\sigma_c(u)$ to the SAME fraction field $L$ by the existing
fraction-field ring-equivalence construction. Its entire value on
each original $b\in B$ is the image of $\sigma_c(u)b$. Uniqueness of
fraction-field maps proves its group law. On the actual chart selected
by $j$, the ratio $g_i/g_j$ is sent to
$u^{a_i(c)-a_j(c)}g_i/g_j$. Hence the actual fraction-adjoin chart
$D_j$ is stable. Its integral closure $E_j$ is stable as well: map the
original monic integral equation using the restricted chart
automorphism and the exact fraction-field square. Applying the inverse
gives a genuine automorphism of the SAME $E_j$.

By (2), this action preserves the actual principal ideal $g_jE_j$.
The group $k^\times$ is rootable in every positive integer degree:
an actual root in the algebraically closed field is nonzero because
its positive power is the supplied unit. Apply the preceding lemma
to the original $E_j$ and $g_jE_j$. It fixes EVERY original minimal
prime, rather than merely permuting the finite family. In particular
it fixes the actual prime in $x$.

The defining original fraction formula for $V_x=(E_j)_P$ now gives
its stability: transform the numerator and denominator of each
$b/s$, with $s\notin P$. The transformed denominator is still outside
the SAME $P$. Scaling by $u^{-1}$ proves the inverse. This constructs
a ring automorphism of the literal original $V_x$, commuting with the
original inclusion from $B$ and the just constructed automorphism of
$B$. The existing DVR automorphism theorem preserves its normalized
order, including zero. Consequently, for EVERY original $H\in B$,
\[
 \eta_x(H)\mid\eta_x(\sigma_c(u)H).
 \tag{3}
\]
This proves the actual threshold stability; it is not supplied as a
certificate.

Fix $F,n,a$ as in (1). The EXISTING parameter-component theorem,
applied to the given actual $t$-monomial family, gives
\[
 \eta_x(F)\mid\eta_x(t^nF_n).
 \tag{4}
\]
In the ORIGINAL ambient polynomial algebra $K[Z_1,\ldots,Z_r]$,
define the $k$-submodule
\[
 N=\{p\in S:\ \eta_x(F)\mid\eta_x(t^np)\}.
\]
It is a genuine submodule of that ambient algebra: membership in $S$
and divisibility are each closed under zero, addition and the original
$k$ scalar action; the polynomial monomial map is additive and
$k$-linear. Formula (4) says that the actual original coefficient
$F_n$ belongs to $N$.

If $p\in N$, let $H=t^np$. Apply the actual automorphism of $V_x$
to the ORIGINAL divisibility $\eta_x(F)\mid\eta_x(H)$. This gives
$\sigma_c(u)(\eta_x(F))\mid\sigma_c(u)(\eta_x(H))$ in that SAME ring.
Formula (3) applied to $F$ gives
$\eta_x(F)\mid\sigma_c(u)(\eta_x(F))$. Transitivity and the exact
coefficient square therefore give
$\eta_x(F)\mid\eta_x(\sigma_c(u)H)$.
Since $\sigma_c(u)(t^np)=t^n\sigma_c(u)p$, this is exactly the
required divisibility of the scaled monomial. Thus $N$ is stable under
every nonzero scaling of each original variable. The EXISTING
multigraded Vandermonde extraction theorem applies to this actual
submodule, over the infinite field $k$. It puts the actual ambient
component $\operatorname{monomial}(a,(F_n)_a)$ in $N$.
Its literal original subtype is exactly
$\operatorname{homComp}_a(F_n)$. This is (1).
All finite supports, zero coefficients and $r=0$ are covered by the
same original submodule and component identities. No new
Vandermonde inversion or analytic point construction is used.
\end{proof}

\begin{remark}
The divisibility theorem concerns EVERY competing ring of the SAME
original cone flag, not only the prescribed divisor. For a jointly
homogeneous cone flag supplied by an original divisorial point, it
provides the missing component comparison in all its competing rings.
The existing original parameter-positivity theorem can then be used
when that flag contains a power of $t$. Constructing the resulting
finite family as actual normalized multigraded analytic points, with
their entire original Gauss identities, is still a separate step.

A native polynomial chart boundary arising from a finite
Fubini--Study family in ONE fixed integer section class already has
other available constructions: simultaneous unit linear coordinates,
whole original Rees-ring transport across actual chart overlaps,
actual common homogeneous representatives, integer-class denominator
units, and degree raising preserving every value on a finite set.
Those constructions do not assert component divisibility in arbitrary
competing cone rings, and the present theorem does not reconstruct
any of them.

To use those native constructions for an arbitrary original cone
flag, one must still derive the actual cone-to-native boundary
comparison, retain and deduplicate the WHOLE original valuation rings,
and prove that all original normalized Gauss values and strict
separator thresholds correspond. Quantitative transfer to a prescribed
polarized degree is also separate. No finite-FS boundary, homogeneous
separator, common normalization, section-surjectivity or degree bound
is supplied as an assumption or a conclusion here. In particular no
global PL barrier, nonnegligibility or envelope-criterion converse is
asserted.
\end{remark}
